% Building a faceted grain-boundary mesostate: an outline.
%
% A merged, math-free summary of gbShiftsEnumeration.tex (which pairs exist)
% and plasticDisplacement.tex (how the surface and its displacement field are
% built).  Consult those two for the equations.
%
% Self-contained: compiles with pdflatex.

\documentclass[11pt]{article}
\usepackage[margin=1.15in]{geometry}
\usepackage{enumitem}
\setlist[itemize]{itemsep=2pt, topsep=4pt, parsep=0pt}
\setlist[enumerate]{itemsep=3pt, topsep=4pt, parsep=0pt}

\title{Building a faceted grain-boundary mesostate\\
       \large an outline of the procedure}
\date{}

\begin{document}
\maketitle

\section*{The idea}

A grain boundary is described by choosing, at a set of sites along it, how far
the two grains are translated relative to one another. Each choice is a
\emph{translation--shift pair}: a relative translation of the two lattices, and
the position along the boundary at which that translation is applied. Engaging
a set of such pairs defines a \emph{mesostate}.

The two grains are then bounded by triangulated surfaces --- the \emph{facets}
--- carrying those translations as a displacement jump. Deforming each grain by
the field that surface generates brings the two boundaries into coincidence and
glues them into one continuous interface, generally stepped or faceted rather
than flat.

The procedure below runs in five stages: enumerate the available pairs, choose
which of them to engage, build the facets, evaluate the displacement, and write
out the atoms.

\section*{Stage 1 --- Enumerating the pairs}

Translations are drawn from the lattice that contains both grains; shifts are
positions on the coincidence lattice. Every candidate translation is combined
with every coincidence site in the boundary cell, and the resulting pair is kept
if its shift lies within an allowed layer about the boundary plane.

Two things are worth understanding about this stage.

\begin{itemize}
  \item \textbf{Every translation is admissible.} Whatever direction a
        translation points in, the two nodes it produces always land on their
        respective lattices. Widening the search therefore cannot produce
        invalid geometry --- it can only produce more pairs.
  \item \textbf{Pairing with every coincidence site is essential.} If the shift
        depended on the translation alone, all the sites generated by one
        translation would collapse onto a single point. A flat boundary, on
        which many sites share one translation, would then be unreachable.
\end{itemize}

Two searches are available.

\begin{itemize}
  \item \textbf{Flat.} Translations are taken from a box whose edges follow the
        directions of the coincidence cell. Because the domain inherits that
        structure, the surviving shifts lie in the boundary plane, and only flat
        boundaries can be built.
  \item \textbf{Full.} That alignment is dropped. Translations are taken from a
        ball, optionally thinned to a slab about the boundary, and enumerated
        directly. Shifts may now retain a component off the boundary plane,
        which is what makes faceted boundaries reachable.
\end{itemize}

The full search needs one extra step. The natural cell in which shifts are
reduced is sheared, so a reduced shift can still sit outside the box that was
asked for. A second reduction folds it back using the two in-plane periods.
This cannot disturb the layer test, because those periods lie in the boundary
plane and so leave the off-plane component untouched.

\section*{Stage 2 --- Choosing which pairs to engage}

Not every combination of pairs is a valid mesostate. Each pair places one node
on each grain, and two engaged pairs may not place nodes on the same site of
either grain --- doing so puts two nodes at one point.

\begin{itemize}
  \item Both sites of a pair depend on the pair alone, so both can be worked out
        before anything is built.
  \item The restriction is pairwise: a set is valid exactly when no two of its
        members collide.
  \item \emph{Both} grains matter. Two pairs can share a site on one grain while
        differing on the other, so checking only one of them leaves half the
        collisions undetected.
\end{itemize}

Enumerating over sets that satisfy this from the outset, rather than building
combinations and discarding the invalid ones, is what makes the sweep
affordable: it also caps the size of a mesostate at the smaller of the two site
counts, which is usually a much lower ceiling than the number of pairs.

\section*{Stage 3 --- Building the facets}

\begin{enumerate}
  \item \textbf{Triangulate the deformed nodes.} The nodes are projected onto
        the boundary plane and triangulated periodically, so the surface wraps
        around the boundary cell rather than having edges. Triangles may span
        the cell boundary; a corner records which periodic copy of a node it
        refers to, so no node has to be duplicated.
  \item \textbf{Share one connectivity between both grains.} Corresponding nodes
        of the two grains deform onto the same point, so giving both surfaces
        the same triangles makes their deformed forms identical --- one
        continuous interface. Triangulating each grain separately does not
        achieve this: their undeformed node sets differ, so they would agree only
        at the nodes and not in between.
  \item \textbf{Orient the two surfaces oppositely.} Each grain must lie on the
        positive side of its own surface, or the sign of its displacement comes
        out backwards.
  \item \textbf{Fix orientation combinatorially, not geometrically.} Some
        triangles are flat in projection and have no orientation of their own, so
        consistency is propagated across shared edges; the geometry is consulted
        once, on the largest triangle, only to decide which way the surface as a
        whole should face.
  \item \textbf{Subdivide for the integration.} The triangles join the nodes, so
        they are as large as the node spacing --- far too coarse to represent a
        displacement that varies from node to node. Each is split into smaller
        panels, across which the displacement is interpolated. This changes no
        geometry: the triangles are flat, so subdividing them only refines the
        integration.
\end{enumerate}

\section*{Stage 4 --- Evaluating the displacement at a point}

The surface acts as a sheet across which the displacement jumps. The field it
generates at a point is the jump weighted by the solid angle the sheet subtends
there. In practice:

\begin{enumerate}
  \item \textbf{Decide which grain the point belongs to.} A point on the far side
        of the boundary belongs to neither grain and is displaced by nothing.
        Everything that follows uses only that grain's surface.
  \item \textbf{Check whether the point is a node.} If it is, return that node's
        own prescribed displacement and stop. This is necessary, not an
        optimisation: the field is genuinely undefined on the surface, where the
        solid angle jumps from one extreme to the other. The nodes are precisely
        the atoms the boundary brings into coincidence, and their displacements
        are already known.
  \item \textbf{Fold the point onto the boundary cell.} The field is periodic, so
        this changes nothing --- but it makes the truncated sum below periodic
        too, which it otherwise would not be.
  \item \textbf{Sum over periodic copies of the surface.} A single finite patch
        subtends a solid angle that dies away with distance instead of tending to
        that of an infinite sheet, so the copies are not a refinement but a
        requirement. Each panel's contribution is computed in closed form, which
        is exact at any distance; a numerical rule fails badly close to the
        surface. Only the nearest copies use the subdivided panels --- further
        out the field varies too little across a triangle for the subdivision to
        matter.
  \item \textbf{Take the sheet's total solid angle from the geometry.} It is
        known exactly for an infinite sheet, whatever its shape, so only which
        side the point lies on has to be determined. This is settled by counting
        how often a ray leaving the point crosses the surface. A
        nearest-surface test is not reliable: inside a corrugation the nearest
        piece of surface can face the wrong way.
  \item \textbf{Close off the remainder analytically.} The copies summed
        explicitly leave a tail that dies away only slowly. Distant copies see
        only the average jump, and the total is known, so the tail is replaced by
        the average jump times the solid angle still unaccounted for. For a
        uniform jump this is exact however few copies are summed, which is the
        case the whole construction has to reproduce: a grain translated bodily
        must come out translated by exactly that amount.
\end{enumerate}

\section*{Stage 5 --- Writing out the atoms}

\begin{enumerate}
  \item \textbf{Cut each lattice back to its own grain.} Both lattices fill the
        whole box, so atoms lying beyond the boundary are discarded. Without
        this the grains interpenetrate and the box holds roughly twice the atoms
        it should.
  \item \textbf{Keep atoms lying on the surface.} The side test has no third
        answer, so an atom exactly on the boundary would be assigned
        arbitrarily. Those atoms are the nodes, and they belong to the grain.
  \item \textbf{Displace every remaining atom} by the field of the surface
        bounding its grain.
  \item \textbf{Wrap positions back into the box in the periodic directions.} A
        displacement can carry an atom across a period boundary, and the file is
        read by a simulation that expects its atoms inside the box it declares.
  \item \textbf{Extend the box in the remaining direction instead of wrapping
        it.} That direction stacks the two grains, so folding an atom across it
        would drop it into the other grain. The box is grown symmetrically,
        only as far as needed, and not at all when nothing lies outside.
\end{enumerate}

\section*{The parameters, and what each controls}

\begin{itemize}
  \item \textbf{Translation radius} --- how far the grains may be translated
        relative to one another. Together with the slab thickness this sets how
        many pairs exist, and so the size of the whole problem.
  \item \textbf{Slab thickness} --- restricts translations to those nearly
        parallel to the boundary. Inert unless smaller than the radius.
  \item \textbf{Layer thickness} --- how far off the boundary plane a shift may
        sit. Reducing it to nothing recovers a flat-only search whichever search
        is selected.
  \item \textbf{Maximum engaged pairs} --- how many pairs one mesostate may
        engage. The clash rule already imposes a ceiling; asking for more than
        that yields nothing.
  \item \textbf{Panel subdivision} --- how finely each triangle is split for the
        integration. Governs how well a displacement varying between neighbouring
        nodes is resolved.
  \item \textbf{Periodic copies} --- how many copies of the surface are summed
        explicitly before the analytic closure takes over. Governs the far field
        only, and only its variation about the average.
\end{itemize}

The last two are balanced against each other: refining the panels while summing
too few copies, or the reverse, leaves one error dominating and wastes the
effort spent on the other.

\end{document}
